El baloncesto es un juego de dos equipos de cinco jugadores del cual se compite por obtener la mayor cantidad de puntos mediante acciones ofensivas como defensivas.

La Women's National Basquetball Association (WNBA), es una liga profesional que consta de 32 partidos, compuesta por 12 equipos cada equipo consta de una plantilla de 12 jugadoras, de estos equipos existen conferencias conocida la Conferencia Este y la Conferencia Oeste entre estos se encuentran:
\begin{itemize}[label=\tiny$\blacksquare$]
    \item Atlanta Dream (Eastern Conference)
    \item Chicago Sky (Eastern Conference)
    \item Connecticut Sun (Eastern Conference)
    \item Indiana Fever (Eastern Conference)
    \item Los Angeles Sparks (Western Conference)
    \item New York Liberty (Eastern Conference)
    \item Minnesota Lynx (Western Conference)
    \item Phoenix Mercury (Western Conference)
    \item San Antonio Stars (Western Conference)
    \item Seattle Storm (Western Conference)
    \item Tulsa Shock (Western Conference)
    \item Washington Mystics (Eastern Conference)
\end{itemize}

Una característica del baloncesto es la posibilidad de modificar las combinaciones de las jugadoras mediante sustituciones. Cada uno de estos cambios genera un cambio en las responsabilidades ofensivas y defensivas, como las relaciones e interacciones entre las jugadoras.
\cite{wnba_faq}
\subsubsection{Posiciones y funciones de las jugadoras}
Tradicionalmente las posiciones de las jugadoras se dividen en:
\begin{itemize}
    \item Point Guard: base o armadora
    \item Shooting Guard: escolta
    \item Small Forward: alera
    \item Power Forward: ala-pivot
    \item Center: pivot
\end{itemize}
Estas posiciones se les asocian con ciertas responsabilidades, como la conducción de el balón, generación de oportunidades de anotación, lanzamiento y protección de la canasta.
\cite{achamps_basketball_positions}
\subsubsection{Alineaciones y sustituciones en el baloncesto}
Cada alineación esta compuesta de 5 jugadoras, que representan al equipo dentro de la cancha. Estas alineaciones permiten considerar elementos como la capacidad de anotación, creación de oportunidades ofensivas, capacidad defensiva, entre otros.

A lo largo de un partido se permite realizar la cantidad de sustituciones deseadas por el entrenador, siempre que el juego este detenido; esto permite alterar las estadisticas de cada quinteto, sin que necesariamente cambie la estrategia general del juego. \cite{wnba_faq}

\subsubsection{Rendimiento individual y colectivo}
El rendimiento individual se refiere a las contribuciones de una jugadora durante un encuentro, con la medicion de estadisticas como: puntos anotados, asistencias, rebotes, robos, bloqueos y perdidas de balon.

Por otra parte, el rendimiento colectivo considera el comportamiento de una combinación de jugadoras cuando participan simultáneamente. Este rendimiento comprende los resultados obtenidos durante las posesiones ofensivas y defensivas del equipo.

Ambas dimensiones se encuentran relacionadas, pero no son equivalentes. Una jugadora puede registrar estadísticas individuales favorables sin que ello implique necesariamente que todas las alineaciones en las que participa presenten el mismo rendimiento. Por esta razón, el análisis colectivo requiere estudiar las combinaciones y los contextos en que se producen las contribuciones individuales.